\subsection{Обратное преобразование Лапласа}

Для исследования переходных процессов воспользуемся найденной ранее системой уравнений движения. Вектор искомых величин определяется выражением

\begin{equation}
    X(p, k_{\text{п}}, k_{\text{и}})
    =
    A^{-1}(p, k_{\text{п}}, k_{\text{и}})
    \cdot
    h(k_{\text{п}}, k_{\text{и}})
    \cdot
    \frac{\pi}{p},
\end{equation}

где $A(p,k_{\text{п}},k_{\text{и}})$ --- матрица системы, а
$h(k_{\text{п}},k_{\text{и}})$ --- вектор входного воздействия.

Для получения временных зависимостей применим обратное преобразование Лапласа:

\begin{equation}
    x(t)
    =
    \mathcal{L}^{-1}\{X(p)\}.
\end{equation}

Исследование переходных процессов проводилось в среде Wolfram Mathematica. Для параметров

\[
    k_{\text{п}} = 0.1,
    \qquad
    k_{\text{и}} = 0.1
\]

были получены графики переходных характеристик системы.

Вектор состояния имеет компоненты:

\[
\begin{aligned}
    x_1(t) &= \varphi_{\text{пр}}(t), \\
    x_2(t) &= \varphi_{\text{м}}(t), \\
    x_3(t) &= M_{\text{дв}}(t), \\
    x_4(t) &= u(t).
\end{aligned}
\]

На рисунках~\ref{fig:pdf/phi_pr_01}--\ref{fig:pdf/u_01} представлены переходные процессы при $k_{\text{п}}=0.1$, $k_{\text{и}}=0.1$.

\screenshot[0.75]{pdf/phi_pr_01}{Переходный процесс $\varphi_{\text{рпр}}(t)$ при $k_{\text{п}}=0.1$, $k_{\text{и}}=0.1$}

\screenshot[0.75]{pdf/phi_m_01}{Переходный процесс $\varphi_{\text{м}}(t)$ при $k_{\text{п}}=0.1$, $k_{\text{и}}=0.1$}

\screenshot[0.75]{pdf/mdv_01}{Переходный процесс $M_{\text{дпр}}(t)$ при $k_{\text{п}}=0.1$, $k_{\text{и}}=0.1$}

\screenshot[0.75]{pdf/u_01}{Переходный процесс $u(t)$ при $k_{\text{п}}=0.1$, $k_{\text{и}}=0.1$}

По полученным графикам видно, что углы $\varphi_{\text{рпр}}(t)$ и $\varphi_{\text{м}}(t)$ стремятся к заданному значению $\pi$. Момент двигателя $M_{\text{дпр}}(t)$ и управляющий сигнал $u(t)$ после начального выброса стремятся к нулю. Следовательно, при выбранных коэффициентах регулятора система является устойчивой.
